% ID: SEL-D0005-OA-E02
% ESTADO: validada. S006, 2026-11-07; tres pasadas del mismo autor y revisión matemática, didáctica y visual.
\textbf{a) Condiciones de extremo y pendiente.}\par\nopagebreak[4]
La derivada es $f'(x)=3x^2+2ax+b$. El extremo relativo en 1 exige $f'(1)=0$; la pendiente de la tangente en 0 es $f'(0)=-1$. Así,
\[
3+2a+b=0,\qquad b=-1\quad\Longrightarrow\quad\boxed{a=-1,\ b=-1}.
\]
Además, $f''(1)=6+2a=4>0$, por lo que el punto estacionario es efectivamente un mínimo relativo, no sólo un candidato.

\textbf{b) Monotonía para esos parámetros.}\par\nopagebreak[4]
El dominio es $\mathbb R$ y
\[
f(x)=x^3-x^2-x+1,\qquad f'(x)=3x^2-2x-1=(3x+1)(x-1).
\]
Los ceros de la derivada son $-1/3$ y 1. Su signo es positivo en $(-\infty,-1/3)$, negativo en $(-1/3,1)$ y positivo en $(1,\infty)$. Por tanto, la función crece en el primer y tercer intervalo y decrece en el segundo.

Hay máximo relativo en $x=-1/3$, con $f(-1/3)=32/27$, y mínimo relativo en $x=1$, con $f(1)=0$. No son extremos absolutos en $\mathbb R$, pues la cúbica no está acotada.

\textbf{Curvatura.}\par\nopagebreak[4]
\[
f''(x)=6x-2,\qquad f''(x)=0\iff x=1/3.
\]
Para $x<1/3$ la segunda derivada es negativa (curvatura hacia abajo); para $x>1/3$ es positiva (hacia arriba). El signo cambia y la función es continua: hay inflexión en $(1/3,16/27)$.

Comprobación: $f'(-1/3)=f'(1)=0$, $f''(-1/3)=-4<0$ y $f''(1)=4>0$, coherentes con los extremos obtenidos. El signo de $f'$ estudia monotonía; el de $f''$, curvatura.
